% SPDX-License-Identifier: CC-BY-4.0
\section{The deduction}

\begin{proposition}\label{prop:answer}
For every pair of positive integers \(n,k\), every polynomial
\(P\in\F_2[x_1,\ldots,x_n]\) satisfying~\eqref{eq:hypotheses}
obeys
\[
 \deg P\ge 2k-\left\lfloor\frac{k}{2^{n-1}}\right\rfloor.
\]
\end{proposition}

\begin{proof}
The condition at the origin excludes the zero polynomial. Put
\(d=\deg P\) and apply Lemma~\ref{lem:msz} with \(S=\F_2\).
Since multiplicities are nonnegative and \(\F_2^n\) has \(2^n-1\)
nonzero points,
\[
 k(2^n-1)
 \le \sum_{a\in\F_2^n}\mult_a(P)
 \le d\,2^{n-1}.
\]
Dividing by \(2^{n-1}\) and using integrality of \(d\) gives
\[
 d\ge\left\lceil 2k-\frac{k}{2^{n-1}}\right\rceil
   =2k-\left\lfloor\frac{k}{2^{n-1}}\right\rfloor.
\]
\end{proof}

Proposition~\ref{prop:answer} answers the published
Question~4.3 of~\cite{BBDM23} without the assumption \(k\ge2^{n-2}\).
The upper bound on the origin's multiplicity is used here only to
ensure that \(P\ne0\). In fact, the same argument applies to any nonzero
polynomial with the required multiplicities at the other points.

The contribution of this note is this direct connection between the
published question and the existing lemma.

\paragraph{Assistance disclosure.}
The author directed the investigation. The literature comparison,
identification of the implication, and drafting were carried out with
substantial assistance from GPT-6 Astra in Codex. This version has not
received independent mathematical review and is not formally verified.
See the \href{https://github.com/kbr-/math-research/blob/main/publications/binary-polynomial-multiplicity/DEVELOPMENT_HISTORY.md}{development history in the Noemesis research repository}.
